\section{Discussion and Limitations}
The central implication is that dense aerial sidelink should jointly control \emph{who communicates with whom} and \emph{how resources are reused}. Resource isolation and spatial selectivity are effective under the long-link baseline, but topology-aware pairing can improve the shared-resource operating point even more dramatically by shortening desired links. This motivates clustering, neighbor selection, relay selection, and reuse coordination as coupled design variables rather than treating swarm size as the sole stressor.

The non-monotonic $R=2$ versus $R=4$ goodput result is also instructive. Orthogonalization is not free: fewer co-channel interferers increase SINR, but narrower resource partitions reduce TBS. A useful scheduler should therefore optimize interference reduction and payload bandwidth jointly rather than maximizing the number of orthogonal groups. Similarly, directionality should be interpreted as residual-interference control after topology and resource reuse have already shaped the interference field.

The topology result changes how a "maximum swarm size" should be interpreted. A single number is not meaningful without a desired-link rule, area/density model, traffic activity, and scheduling policy. The operating envelope in Fig.~\ref{fig:envelope} is best viewed as a controlled engineering map for one baseline topology, while Fig.~\ref{fig:topology} demonstrates how strongly that map can move when local communication is enforced.

\subsection{Threats to validity and generalization}
The main external-validity limitation is geometry. Increasing $N$ in a fixed $1000\times1000$~m area increases spatial density, while sequential pairing leaves desired-link length roughly unchanged. Nearest-neighbor pairing reveals the opposite geometric effect: as density rises, local desired links become shorter. A different scaling law would be expected if physical area grew with $N$, if communication partners were constrained by mission roles, or if only a fraction of links were simultaneously active. The reported curves therefore characterize two explicit topology families rather than all possible swarm organizations.

The propagation model is large-scale and equal-altitude. It is anchored in measured A2A data, but the simulated positions are synthetic and the study omits shadowing correlation, small-scale fast fading, Doppler, blockage, antenna-body interaction, and trajectory-dependent channel memory. These effects can widen the instantaneous SINR distribution and may change outage tails even when the mean large-scale trend is preserved. Future evaluation should therefore repeat the matched-seed design with standards-based fading, mobility traces, and measured multi-UAV RF data.

The PHY/MAC abstraction creates a second validity boundary. 5G-LENA numerical BLER curves are link-level simulation evidence, not standardized BLER requirements. Transport-block BLER is reconstructed from sourced code-block curves under an independent-code-block approximation, and the adaptive MCS controller assumes immediate knowledge of the modeled SINR. The fixed-MCS control shows that this assumption is performance-relevant. Likewise, the weighted conflict graph is a transparent geometry proxy, not NR sidelink Mode-2 sensing, and the directional parameter $G$ is not a beam codebook, array pattern, tracking loop, or measured antenna gain. Consequently, absolute throughput values should not be interpreted as deployment guarantees.

\subsection{Design implications}
Within these boundaries, several design implications are stable. First, topology control should precede any claim about swarm capacity: local neighbor communication can offset a substantial part of the interference penalty associated with larger $N$. Second, orthogonalization should be applied selectively because the SINR benefit competes with reduced PRBs and TBS; the $R=2$ versus $R=4$ result is a concrete example where more isolation does not immediately increase goodput. Third, geometry-aware reuse provides additional gain beyond random partitioning, suggesting that even lightweight interference-aware grouping can be valuable before introducing more complex schedulers. Fourth, spatial selectivity and link adaptation should be treated as complementary controls, not substitutes for topology and resource design.

A practical implementation could therefore operate hierarchically: maintain short communication neighborhoods where the mission permits; estimate local conflict or sensing metrics; choose a reuse level that balances interference against bandwidth; then apply spatial selectivity and MCS adaptation to residual interference. The present paper does not claim that the specific conflict graph or idealized $G$ abstraction is the final protocol. Their purpose is to expose the value of each control dimension under reproducible conditions so that a standards-complete scheduler can later be evaluated against the same operating-region methodology.

The study is intentionally system-level. It does not implement a bit-accurate sidelink PHY, normative Mode-1/Mode-2 sensing/scheduling, full fast fading, full MIMO/beam management, or measured multi-UAV RF interference/PDR. BLER curves are link-level simulation evidence, not 3GPP-standard BLER tables. The conflict-graph allocator, adaptive MCS rule, pairing rules, and $G$ sweep are research abstractions. The operating-envelope thresholds are engineering-policy choices. These limitations bound the claims but also make the experiment auditable: each mechanism can be changed independently and re-evaluated under matched seeds.
